%==============================================================================
% 3. MATERIALES Y MÉTODOS
%
% Pauta: "Indicar el material utilizado, señalando su procedencia y los métodos
% con la bibliografía correspondiente. Los métodos nuevos o modificaciones
% deben describirse en detalle."
%==============================================================================

\chapter{Materiales y Métodos}

\section{Área de estudio}

Describe el área de estudio, sus coordenadas y sus características
oceanográficas relevantes.

% ---------------------------------------------------------------------------
% EJEMPLO DE FIGURA
%
% Pauta: la numeración, el título y la leyenda van BAJO el cuerpo de la figura,
% con números arábigos y en castellano. La figura debe ser autoexplicativa e
% indicar la fuente si no es original.
%
% Atajo de la plantilla:
%   \figura[ancho]{archivo}{Título. Leyenda…}{etiqueta}
% ---------------------------------------------------------------------------
%
% \figura[0.8\linewidth]{area_estudio.png}
%   {Área de estudio. Batimetría de la plataforma continental frente a
%    Valdivia (39,8\textdegree\,S) y ubicación de las estaciones de muestreo
%    (puntos negros). Fuente: elaboración propia a partir de GEBCO 2023.}
%   {area-estudio}
%
% Y en el texto se cita así: «\dots como se muestra en la \reffig{area-estudio}»
% o bien «\dots (\reffig{area-estudio})».

\section{Origen de los datos}

Indica la procedencia de cada conjunto de datos: sensor o instrumento, período,
resolución espacial y temporal, y la referencia correspondiente.

% ---------------------------------------------------------------------------
% EJEMPLO DE TABLA
%
% Pauta: las Tablas se numeran con números ROMANOS EN MAYÚSCULAS y la
% numeración, junto con el título, se incluye SOBRE el cuerpo de la Tabla.
%
% Atajo de la plantilla:
%   \begin{tablabima}{Título de la tabla}{etiqueta} … \end{tablabima}
% ---------------------------------------------------------------------------

\begin{tablabima}{Conjuntos de datos utilizados, con su procedencia,
  resolución y período de cobertura.}{datos}
  \begin{tabular}{llll}
    \toprule
    \textbf{Producto} & \textbf{Variable} & \textbf{Resolución} & \textbf{Período} \\
    \midrule
    ERA5   & Viento a 10 m        & 0,25\textdegree{} / 1 h & 1990--2024 \\
    MODIS  & Temperatura sup.\ del mar & 4 km / 1 d        & 2002--2024 \\
    \bottomrule
  \end{tabular}
\end{tablabima}

\noindent
En el texto, la tabla se cita igual que una figura:
«\dots los productos empleados se resumen en la \reftab{datos}».

\section{Procesamiento y análisis}

Describe en detalle los métodos aplicados, citando la bibliografía
correspondiente \parencite{Conover1980}. Los métodos nuevos, o las
modificaciones que hayas introducido a métodos existentes, deben quedar
descritos con el suficiente detalle como para que otra persona pueda
reproducirlos; el procedimiento sigue el descrito por \textcite{Villegas1988}
y \textcite{Permitin1970}.

Las unidades se escriben con \verb|siunitx|, por ejemplo
\verb|\qty{12.4}{\degreeCelsius}| $\rightarrow$ \qty{12.4}{\degreeCelsius} o
\verb|\qty{0.25}{\meter\per\second}| $\rightarrow$ \qty{0.25}{\meter\per\second}.
